\documentclass[11pt,a4paper]{article}
\usepackage[margin=2.2cm]{geometry}
\usepackage{listings}
\usepackage{xcolor}
\usepackage[colorlinks=true,linkcolor=blue]{hyperref}

\lstset{
  basicstyle=\ttfamily\small,
  backgroundcolor=\color{black!5},
  frame=single,
  rulecolor=\color{black!20},
  breaklines=true,
  columns=fullflexible,
  keepspaces=true
}
\setlength{\parindent}{0pt}
\setlength{\parskip}{6pt}

\title{Running the Lidar and ZED 2i}
\author{Rover Drivers Quick Guide}
\date{}

\begin{document}
\maketitle

\section{Before you start}

Build the workspace once (from the repo root) and source it in every new terminal:

\begin{lstlisting}[language=bash]
rosdep install --from-paths src --ignore-src -r -y
colcon build --symlink-install --cmake-args -DCMAKE_BUILD_TYPE=Release
source install/setup.bash
\end{lstlisting}

The ZED needs the \textbf{ZED SDK} and \textbf{CUDA} installed on the machine. If the build fails on the ZED part, this is usually why.

\section{Unitree 4D Lidar L2}

\subsection{Connect it}

The lidar talks over Ethernet (UDP) by default.

\begin{itemize}
  \item Plug the lidar into the computer with the Ethernet cable and power it on.
  \item Lidar IP: \texttt{192.168.1.62}
  \item Set your computer's wired IP to a fixed address: \texttt{192.168.1.2}, netmask \texttt{255.255.255.0}.
  \item Check the link: \texttt{ping 192.168.1.62}
\end{itemize}

\subsection{Run it}

\begin{lstlisting}[language=bash]
ros2 launch unitree_lidar_ros2 launch.py
\end{lstlisting}

This starts the driver and opens RViz with the point cloud.

\subsection{Check it works}

\begin{lstlisting}[language=bash]
ros2 topic list
ros2 topic hz /unilidar/cloud
ros2 topic hz /unilidar/imu
\end{lstlisting}

\begin{tabular}{ll}
  Point cloud & \texttt{/unilidar/cloud} (frame \texttt{unilidar\_lidar}) \\
  IMU         & \texttt{/unilidar/imu} (frame \texttt{unilidar\_imu}) \\
\end{tabular}

\subsection{Using USB/serial instead}

In the launch file (\texttt{unitree\_lidar\_ros2/launch/launch.py}), \texttt{initialize\_type} is \texttt{2} for UDP. Change it to \texttt{1} for serial. The serial port is \texttt{/dev/ttyACM0}.

\section{ZED 2i}

\subsection{Connect it}

\begin{itemize}
  \item Plug the camera into a \textbf{USB 3} port (blue port, USB 3 cable).
  \item Check the camera is seen: \texttt{lsusb | grep -i stereolabs}
\end{itemize}

\subsection{Run it}

\begin{lstlisting}[language=bash]
ros2 launch zed_wrapper zed_camera.launch.py camera_model:=zed2i
\end{lstlisting}

The first start can take a few minutes because the SDK optimizes its neural depth model. After that it starts fast.

\subsection{Check it works}

\begin{lstlisting}[language=bash]
ros2 topic list | grep zed
ros2 topic hz /zed/zed_node/rgb/color/rect/image
\end{lstlisting}

Useful topics:

\begin{tabular}{ll}
  Color image & \texttt{/zed/zed\_node/rgb/color/rect/image} \\
  Depth       & \texttt{/zed/zed\_node/depth/depth\_registered} \\
  Point cloud & \texttt{/zed/zed\_node/point\_cloud/cloud\_registered} \\
  IMU         & \texttt{/zed/zed\_node/imu/data} \\
  Odometry    & \texttt{/zed/zed\_node/odom} \\
\end{tabular}

Look at it in RViz: \texttt{rviz2}, then add the image or point cloud topic.

\section{Run both together}

Open two terminals (source the workspace in each) and run the lidar launch in one, the ZED launch in the other. They do not clash: the lidar uses Ethernet, the ZED uses USB.

\section{Common problems}

\begin{tabular}{p{5cm}p{9cm}}
  \textbf{Problem} & \textbf{Fix} \\ \hline
  No lidar data & Check the computer IP is \texttt{192.168.1.2} and \texttt{ping 192.168.1.62} works. \\
  Lidar serial permission error & \texttt{sudo usermod -aG dialout \$USER}, then log out and in. \\
  ZED not found & Use a USB 3 port and cable. Replug the camera. \\
  ZED build error & ZED SDK or CUDA is missing or the versions do not match. \\
  ZED slow first start & Normal. It is optimizing the model once. \\
  Topics missing & Did you \texttt{source install/setup.bash} in this terminal? \\
\end{tabular}

\end{document}
